\section{Introduction}

% ==========================================
% 1. SECTION 1.1
% ==========================================

\subsection{Predictive Modeling in Low-Data Bioactivity Regimes}

Quantitative structure--activity relationship (QSAR) modeling is constrained when few labeled compounds are available
\citep{winkler2026}. In that regime, simpler and more interpretable approaches, such as multivariate linear regression
or classical tree-based ensembles, often match or exceed the predictive reliability of neural networks
\citep{hadas2023, xiewt2023}. These models carry more bias and less variance, and they are less sensitive to the random
composition of the training sample than overparameterized alternatives, particularly when they are given descriptors
derived from quantum-chemical calculations \citep{xiewt2023}.

Small data sets invite overfitting, and the risk grows as the number of molecular descriptors approaches or exceeds
the number of compounds \citep{hawkins2004, godarzi2013}. One common response is to enlarge the training set by data
augmentation or synthetic sample generation \citep{gerz2026, tyler2025}, which costs computation \citep{tyler2025,
calle2025} and, applied carelessly, obscures what the base algorithm can do on its own. When augmentation precedes
partitioning, the ``augment-then-split'' fallacy, perturbed copies of a test molecule reach the training set; the
partitions are then no longer independent and the reported metrics are inflated \citep{gerz2026, john2025,
kannoth2024}. Augmentation therefore has to stay inside the training folds of a partition fixed beforehand
\citep{gerz2026, john2025}.

Evaluating complex architectures in this regime also requires separating model building from model assessment. Tuning
hyperparameters on a small sample with a large search space biases the selection \citep{cawley2010}, so an architecture
scored on the same data that chose it says little about its behavior on unseen compounds. Double, or nested,
cross-validation is the established way to decouple the two \citep{bauman2014, cawley2010}. Even then, a random
$k$-fold partition measures interpolation among related compounds rather than extrapolation to new chemotypes
\citep{kyrali2022, xiong2020}, and generative deep learning extrapolates poorly in the sparse, discrete space of
Simplified Molecular-Input Line-Entry System (SMILES) strings \citep{yifan2023}.

A model fitted to few compounds can also reach a good score by chance correlation, which an empirical check can
detect. In y-randomization the response is permuted and the model refitted, so that a score close to the one obtained
with the true responses points to chance correlation rather than to a learned relationship \citep{rucker2007}. R\"ucker
et al.\ prescribe more than permuting the response: the whole model-building procedure, descriptor selection included,
should be repeated inside each randomization, and the result compared with models built on random descriptors
\citep{rucker2007}. Where curated descriptors pass these checks, machine learning has been used to rank candidate
compounds, for instance to separate amyloid-$\beta$ aggregation inhibitors from decoys \citep{avantika2025}.

Transfer learning offers a different route around data scarcity, by adapting networks trained on a large source domain
\citep{samuel2023}. Reduced learning rates during fine-tuning \citep{damilola2025} and lighter convolutional
architectures in place of Transformers \citep{riza2025} have been reported to help with class imbalance and selection
bias in low-resource settings.

% ==========================================
% 1. SECTION 1.2
% ==========================================
\subsection{Domain-Engineered Features versus End-to-End Representation Learning}

Feature-based machine learning relies on domain-engineered descriptors and molecular fingerprints \citep{rajesh2026}.
Morgan (ECFP4) fingerprints encode local substructures, while physicochemical and 3D descriptors quantify specific
molecular properties; both are computed deterministically and need no large training corpus \citep{avantika2025,
beatrice2026, xiewt2023}. Their drawback is dimensionality: accumulating handcrafted features produces large and
redundant matrices \citep{godarzi2013, hadas2023}, which on restricted data sets raises computational cost and makes
feature selection itself a source of error \citep{godarzi2013, hadas2023}.

End-to-end representation learning avoids manual feature engineering. Variational autoencoders (VAEs), generative
adversarial networks (GANs) and graph neural networks (GNNs) map molecular topologies into continuous latent spaces
and model nonlinear interactions without a predefined descriptor set \citep{abul2026, rajesh2026, yifan2023}, and
chemical language processing (CLP) applies the same idea to string representations such as SMILES \citep{riza2025}. The
discrete nature of those strings is a weakness of the approach, since altering a single character can produce an
invalid structure \citep{yifan2023}.

How much this flexibility is worth when labels are scarce remains unsettled. Representation learning needs enough
examples to recover what a descriptor encodes by construction, which makes it less effective as the labeled set
shrinks \citep{sotirios2025}. For inhibitors of \textit{Toxoplasma gondii} dihydrofolate reductase (TgDHFR), \citet{karimova2025} addressed this limitation with Gaussian noise-based data augmentation and ensembles
of deep neural networks, raising their baseline $R^2$ from \PaperBaselineRTwo{} to \PaperFinalRTwo{}, and released
both the curated data set and the code that produced it. Their result poses the question this study takes up: on a set
of this size, is synthetic data generation necessary for predictive reliability? Classical ensembles are reported to
match or exceed deep architectures when data are scarce \citep{beatrice2026, hadas2023, rajesh2026, salman2023,
samuel2023, xiewt2023}, although such comparisons are rarely run on identical, leakage-free folds.

Opinion on nonlinear models in limited-data settings is itself shifting \citep{Dalmau2025}. Regularized nonlinear
algorithms combined with structured linear priors reach results comparable to classical methods \citep{Dalmau2025,
salman2023}, and fine-tuning pretrained networks such as ChemBERTa has produced usable predictions from a few dozen
labeled compounds \citep{latefa2025, samuel2023}.

% ==========================================
% 1. SECTION 1.3
% ==========================================
\subsection{Data Leakage in Low-Data Molecular Modeling}

Data leakage, the entry of test-set information into the training process, inflates performance estimates without
producing an error or a warning \citep{brva2026, gerz2026}. In low-data QSAR it arises most often in two places. The
first is preprocessing: when scaling, imputation or variance-based filtering is fitted on the complete data set, the
model is given information about the distribution of the test compounds, the independence of the partitions is lost,
and the resulting estimate is optimistic \citep{correa2025, lewis2023}.

The second is model selection. A validation set that repeatedly guides hyperparameter choices stops being independent
of the model, and highly parameterized architectures such as deep neural networks are the best placed to exploit that
overlap, fitting noise and returning optimistic estimates of the generalization error \citep{bauman2014, sotirios2025}.
Both paths are closed by fitting every transformation and every selection step inside the training partition at hand.
Because the details of the partitioning are often reported incompletely, performance is sometimes attributed to a new
architecture when it comes from leakage instead \citep{rojas2026}.

% ==========================================
% 1. SECTION 1.4
% ==========================================

\subsection{Objectives and Principal Contributions}

This study asks whether, on one small and well-characterized bioactivity set (\NCompounds\ TgDHFR inhibitors), graph networks, 3D descriptors, feature selection and data augmentation improve on classical ensembles trained on fingerprints once the evaluation is free of leakage. We do not assume the answer; we design the evaluation so that either outcome would be detectable. Our contributions are:

\begin{itemize}
	\item \textbf{Leakage-free, paired evaluation:} A reproducible Python pipeline built on a $\OuterKFolds \times \OuterRepeats$ repeated nested cross-validation protocol (\TotalOuterFolds{} outer evaluations) in which every preprocessing and selection step is fitted inside the training partitions, and in which the classical models and the D-MPNN are scored on exactly the same folds.

	\item \textbf{Calibrated statistical inference:} Use of the Nadeau--Bengio corrected $t$-test with Holm adjustment for the comparisons, which accounts for the overlap of training sets across repeated folds, and of an equivalence test (TOST) with a pre-specified margin, so that ``no significant difference'' is distinguished from ``demonstrated equivalence''.

	\item \textbf{Classical ensembles versus a graph network:} A comparison of four classical representations against a Chemprop D-MPNN with default hyperparameters and with a hyperparameter search nested in the same folds, extended to a seed ensemble, a descriptor-augmented network, a feed-forward network on the same fingerprints and a pretrained foundation model, together with the measured computational cost of each.

	\item \textbf{Audit of the original pipeline:} Re-execution of the published code, followed by the removal of one information leak at a time, which quantifies what each leak contributes to the reported $R^2$ and how much of it survives.

	\item \textbf{Replication of the original study's gains:} A test of whether the feature selection and Gaussian-noise augmentation reported by \citet{karimova2025}, which took their $R^2$ from \PaperBaselineRTwo{} to \PaperFinalRTwo{} on a single hold-out split, transfer to a classical learner under nested cross-validation, with and without the leak that the original selection contains.

	\item \textbf{Robustness checks:} Repetition of the benchmark under scaffold-grouped splits and with the extreme-potency records retained, to quantify how much the conclusions depend on the splitting scheme and on the curation.

	\item \textbf{Screening as hypothesis generation:} Application of the final model to FDA-approved drugs, followed by molecular docking and in silico ADMET filters, to produce a short list of candidates. These are computational hypotheses; none was tested experimentally, and selectivity against human DHFR was not assessed.
\end{itemize}
